\documentclass[11pt,letterpaper]{article}
\usepackage[margin=0.88in,headheight=15pt]{geometry}
\usepackage[T1]{fontenc}
\usepackage{lmodern,amsmath,amssymb,amsthm,mathtools,bm}
\usepackage{microtype,booktabs,longtable,array,enumitem,fancyhdr}
\usepackage[dvipsnames]{xcolor}
\usepackage{xurl}
\usepackage{hyperref}
\hypersetup{colorlinks=true,linkcolor=MidnightBlue,citecolor=MidnightBlue,urlcolor=MidnightBlue,
 pdftitle={Quadratic-Spectrum Stieltjes Identities},pdfauthor={Prepared for Vladimir Reshetnikov with OpenAI assistance}}
\urlstyle{tt}
\numberwithin{equation}{section}
\newtheorem{theorem}{Theorem}[section]
\newtheorem{proposition}[theorem]{Proposition}
\newtheorem{corollary}[theorem]{Corollary}
\newtheorem{lemma}[theorem]{Lemma}
\theoremstyle{definition}
\newtheorem{definition}[theorem]{Definition}
\newtheorem{remark}[theorem]{Remark}
\newcommand{\Li}{\operatorname{Li}}
\newcommand{\Rea}{\operatorname{Re}}
\newcommand{\Ima}{\operatorname{Im}}
\newcommand{\arsinh}{\operatorname{arsinh}}
\newcommand{\dd}{\,\mathrm{d}}
\newcommand{\Z}{\mathbb Z}
\newcommand{\N}{\mathbb N}
\newcommand{\C}{\mathbb C}
\newcommand{\Bcal}{\mathcal B}
\newcommand{\Hcal}{\mathcal H}
\newcommand{\qfile}[1]{\nolinkurl{#1}}
\newcommand{\qpin}{\nolinkurl{446dca9bd65c20e409d1d2561d4ffad408103e2a}}
\setlist[enumerate]{itemsep=4pt,topsep=5pt}
\setlength{\emergencystretch}{2em}
\pagestyle{fancy}
\fancyhf{}
\fancyhead[L]{\small Quadratic-spectrum Stieltjes identities}
\fancyhead[R]{\small ProveIt research continuation}
\fancyfoot[C]{\thepage}
\title{\textbf{Quadratic-Spectrum Stieltjes Identities}\\[7pt]
\large Complete reflected jets, harmonic counterterms,\\
finite polylogarithmic reductions, and a polynomial parity criterion}
\author{Prepared for Vladimir Reshetnikov's ProveIt programme\\
\small Mathematical derivations and computational audit with OpenAI assistance}
\date{10 October 2026}
\begin{document}
\maketitle
\thispagestyle{plain}
\begin{abstract}
We develop a quadratic-spectrum continuation of the ProveIt programme in
Hurwitz zeta, generalized Stieltjes constants, harmonic sums, and polylogarithmic
integration. A positive mass parameter separates the pole of a bilateral
quadratic zeta function from an exponentially convergent Fourier--Bessel
part. An odd companion supplies the missing reflection channel. Their
spectral derivatives give, for every nonnegative index, exact formulas for
both $\gamma_m(a)+\gamma_m(1-a)$ and $\gamma_m(a)-\gamma_m(1-a)$, and hence
for each individual generalized Stieltjes constant. The formulas hold for
every positive mass, not just inside a small-mass power-series disk.

The counterterms have explicit complete-Bell-polynomial coefficients in
ordinary and half-shifted generalized harmonic numbers. A finite-head
identity gives their global convergent implementation. Bessel order
integrals turn both reflected families into damped polylogarithm order
integrals; parameter primitives connect them to Hurwitz derivatives at
zero and to a convergent log-Gamma ratio formula.

For polynomially weighted one-sided spectra we prove an exact conjugate
factorization defect and classify its universal vanishing: it is zero for
every positive mass precisely for even polynomial weights. This yields
finite identities between complex-shift Hurwitz derivatives and
polylogarithms at every nonpositive integer spectral point. Those values
form an exact mass-antiderivative ladder. We also give positive-integer
polygamma reductions and an ordinary Gamma-product integral evaluated by
a dilogarithm.

The classical Poisson, Bessel, Gamma, and multiplicative-anomaly ingredients
retain their attribution. The contribution is a proved, explicit synthesis
and its specialized identities, not a claim of global priority. The
article includes exact algebra regressions, independent numerical
diagnostics, a scoped source audit, and further research questions.
\end{abstract}
\medskip
\noindent\textbf{Source pin.} Repository reading is anchored at \qpin.
No repository files were modified. The incoming archive catalogue was read;
the six archive interiors were not inspected in this session.
\clearpage
\begingroup\small\tableofcontents\endgroup
\clearpage
% BEGIN sections/01-context
\section{Research target and precise scope}\label{qspec:sec:context}

The canonical manuscript organizes polylogarithms, Gamma products, Hurwitz
spectral derivatives, and generalized Stieltjes constants within one exact
calculus \cite{repo}. Its current research chapter explicitly distinguishes
an analytic identity from a relation in a formal presentation and from a
small numerical residual. We retain that distinction. In particular, the
Gaussian $S_6$ identity is not proved here, and the rejection of one frozen
$S_8$ vector is not interpreted as rejection of every identity in its
comparison space.

The present continuation uses a different exact mechanism: introduce a
positive parameter $b$ into $x^2+b^2$, evaluate the resulting bilateral
spectrum in two ways, and then remove the deformation by a convergent
difference. The even spectrum sees the sum of reflected Hurwitz functions.
A first-moment spectrum sees their difference. Neither channel alone
recovers an individual Stieltjes constant, but together they do.

\subsection{The results and their logical dependence}
Theorem~\ref{qspec:thm:master} is the complete reflected Stieltjes identity.
Theorem~\ref{qspec:thm:polylog-integrals} replaces its Fourier--Bessel terms
by ordinary integrals of damped polylogarithms. The harmonic coefficient
formula and finite-head identity in Theorems~\ref{qspec:thm:harmonic}
and~\ref{qspec:thm:head} make all spectral derivatives explicit and avoid
an invalid extension of a binomial series outside its convergence disk.

The second main result, Theorem~\ref{qspec:thm:anomaly}, gives a complete
polynomial-weight parity criterion for the conjugate factorization defect.
It is independent of the Stieltjes master identity. Combining it with the
classical half-integer Bessel polynomials gives
Theorem~\ref{qspec:thm:finite}, a finite complex-Hurwitz/polylogarithm identity
at every negative integer. Theorem~\ref{qspec:thm:mass-ladder} identifies the
same family as repeated mass antiderivatives of a logarithm.

An illustrative member is obtained by setting
$z=\exp(-2\pi b+2\pi ia)$:
\begin{align}
&2\Rea\bigl\{\zeta'(-2,a+ib)+\zeta'(-2,1-a+ib)\notag\\[-2pt]
&\hspace{35mm}-2ib\bigl(\zeta'(-1,a+ib)+\zeta'(-1,1-a+ib)\bigr)\bigr\}\notag\\
&\qquad=-\frac{4\pi b^3}{3}
-\frac{2b}{\pi}\Rea\Li_2(z)-\frac1{\pi^2}\Rea\Li_3(z).
\label{qspec:eq:opening-example}
\end{align}
It holds for every $0<a<1$ and $b>0$. There is no unreported tail on the
right. Its proof is part of the all-index theorem, not an integer-relation
search.

\subsection{Attribution and what is not asserted}
The Gaussian Poisson identity, Mellin transforms, half-integer Bessel
functions, Hurwitz parameter differentiation, and Gamma reflection are
classical \cite{theta,bessel,halfK,hurwitz,gamma}. The first-derivative
multiplicative anomaly is also a classical subject; Dowker treats products
of linear factors and polynomial degeneracies \cite{dowker}. The repository
already has an arithmetic spectral-jet continuation that separates local
residue defects from nonlocal values \cite{arithmetic}. Accordingly, we do
not advertise the existence of the anomaly, Poisson continuation, or
negative-polygamma primitives as discoveries.

What is established here is an explicit two-channel Stieltjes construction,
its global harmonic implementation, a conjugate-pair polynomial parity
classification, and the resulting finite identities and primitive ladders.
These statements are derived in full below. Their absence from the portions
inspected is not an exhaustive literature or repository priority check.
No transcendence, numerical period independence, or proof-assistant
formalization is claimed.

\subsection{Conventions and analytic domains}\label{qspec:sec:conventions}
Unless stated otherwise,
\begin{equation}\label{qspec:eq:domain}
0<a<1,\qquad b>0,\qquad Q_b(x)=x^2+b^2,
\qquad x_n=n+a,\quad y_n=n+1-a.
\end{equation}
A one-sided formula needs only $a>0$. The variable $s$ is a spectral order;
$\zeta^{(r)}(s,a)$ and $\zeta'(s,a)$ differentiate that order, not $a$.
We use
\begin{equation}\label{qspec:eq:stieltjes-convention}
\zeta(1+u,a)=\frac1u+
\sum_{m=0}^{\infty}\frac{(-1)^m}{m!}\gamma_m(a)u^m,
\qquad G_m^\pm(a)=\gamma_m(a)\pm\gamma_m(1-a).
\end{equation}
In particular $\gamma_0(a)=-\psi(a)$. Superscripts on $\gamma_m^{(p)}$
will mean parameter derivatives.

All powers of positive $Q_b(x)$ use the real logarithm. For complex shifts
$a\pm ib$ we use the principal logarithm, and the branch of $\log\Gamma$
holomorphic on the right half-plane and real on its positive axis. Thus
$\log(x+ib)+\log(x-ib)=\log Q_b(x)$ for $x>0$, with no branch correction.
The damped polylogarithms have arguments strictly inside the unit disk,
where their defining series is entire in the order. Unit-circle limits
are stated separately.

The complete exponential Bell polynomials are defined by
\begin{equation}\label{qspec:eq:bell}
\exp\left(\sum_{r\ge1}v_r\frac{t^r}{r!}\right)
=\sum_{m\ge0}\Bcal_m(v_1,\ldots,v_m)\frac{t^m}{m!}.
\end{equation}
The convention $\Bcal_0=1$ and the recurrence
\begin{equation}\label{qspec:eq:bell-recurrence}
\Bcal_m=\sum_{r=1}^m\binom{m-1}{r-1}v_r\Bcal_{m-r}
\end{equation}
make every finite coefficient used below directly reproducible. We write
$H_k^{(r)}=\sum_{j=1}^k j^{-r}$ and $H_0^{(r)}=0$.

% END sections/01-context

% BEGIN sections/02-transform
\section{Even and odd quadratic spectra}\label{qspec:sec:transform}

Define the absolutely convergent series
\begin{align}
E(s;a,b)&=\sum_{n\in\Z}\bigl((n+a)^2+b^2\bigr)^{-s},
&&\Rea s>\tfrac12,\label{qspec:eq:E}\\
Y(s;a,b)&=\sum_{n\in\Z}(n+a)\bigl((n+a)^2+b^2\bigr)^{-s},
&&\Rea s>1.\label{qspec:eq:Y}
\end{align}
Their reflected positive-index forms are
\[
E(s)=\sum_{n\ge0}\bigl(Q_b(x_n)^{-s}+Q_b(y_n)^{-s}\bigr),
\qquad
Y(s)=\sum_{n\ge0}\bigl(x_nQ_b(x_n)^{-s}-y_nQ_b(y_n)^{-s}\bigr).
\]
Continuation of the initially absolutely convergent functions, rather than
an unannounced summation prescription for divergent series, is used throughout.

\begin{proposition}[Classical Poisson--Bessel continuation in both channels]
\label{qspec:prop:poisson}
The meromorphic continuation of $E$ is
\begin{align}
E(s;a,b)={}&\frac{\sqrt\pi\,\Gamma(s-\tfrac12)}{\Gamma(s)}b^{1-2s}\notag\\
&+\frac{4\pi^s b^{1/2-s}}{\Gamma(s)}
\sum_{n\ge1}n^{s-1/2}\cos(2\pi na)K_{s-1/2}(2\pi bn).
\label{qspec:eq:poisson}
\end{align}
Its only poles are simple, at $s=\tfrac12-j$, $j\ge0$, with
\begin{equation}\label{qspec:eq:residues}
\operatorname*{Res}_{s=1/2-j} E(s;a,b)
=\frac1{4^j}\binom{2j}{j}b^{2j}.
\end{equation}
It vanishes at every nonpositive integer. The odd companion is entire and
has the expression
\begin{equation}\label{qspec:eq:odd-poisson}
Y(s;a,b)=\frac{4\pi^s b^{3/2-s}}{\Gamma(s)}
\sum_{n\ge1}n^{s-1/2}\sin(2\pi na)K_{s-3/2}(2\pi bn).
\end{equation}
All displayed Fourier series and all their spectral derivatives converge
locally uniformly for $s\in\C$, $b>0$, and real $a$.
\end{proposition}
\begin{proof}
Mellin integration in the absolute-convergence half-plane gives
\[
\Gamma(s)E(s;a,b)
=\int_0^\infty t^{s-1}e^{-b^2t}
\sum_{n\in\Z}e^{-t(n+a)^2}\dd t.
\]
Apply the Gaussian Poisson identity
\[
\sum_{n\in\Z}e^{-t(n+a)^2}
=\sqrt{\pi/t}\sum_{k\in\Z}e^{-\pi^2k^2/t}e^{2\pi ika}.
\]
This identity follows from the Fourier transform of a Gaussian, or from
the classical theta transformation \cite{theta}. The zero Fourier mode
is $\sqrt\pi\Gamma(s-\tfrac12)b^{1-2s}$. For $k\ne0$, use the
Bessel integral \cite{bessel}
\[
\int_0^\infty t^{s-3/2}e^{-b^2t-\pi^2k^2/t}\dd t
=2(\pi|k|/b)^{s-1/2}K_{s-1/2}(2\pi b|k|).
\]
Combining the two signs of $k$ proves \eqref{qspec:eq:poisson}.
Interchange of the nonzero modes and the integral is justified by the
same integrals with absolute values; their resulting Bessel series has
exponential decay.

For explicit normal convergence, use
\begin{equation}\label{qspec:eq:K-cosh}
K_\nu(u)=\int_0^\infty e^{-u\cosh t}\cosh(\nu t)\dd t,
\qquad u>0.
\end{equation}
On a compact set of $\nu$ the integrand, and any fixed number of order
derivatives, are bounded by $e^{-u\cosh t}e^{A t}(1+t)^M$ for fixed
$A,M$. For $u\ge u_0>0$ the factor $e^{-u\cosh t}$ supplies an
exponential factor in $u$ and an integrable majorant in $t$. Multiplying
by any fixed power of $n$ or $\log n$ does not destroy convergence when
$u=2\pi bn$ and $b$ stays bounded away from zero. Parameter derivatives
in $a$ introduce only powers of $n$ and have the same property.

The nonzero-mode part is therefore entire in $s$. The poles and residues
come only from $\Gamma(s-\tfrac12)$; the half-integer Gamma values give
\eqref{qspec:eq:residues}. At $s=0,-1,-2,\ldots$ the bracket before
multiplication by $1/\Gamma(s)$ is regular, so $E(s)=0$.

Initially for $\Rea s>\tfrac32$, termwise differentiation gives
\begin{equation}\label{qspec:eq:odd-differentiation}
Y(s;a,b)=-\frac1{2(s-1)}\partial_a E(s-1;a,b).
\end{equation}
The zero-mode term in \eqref{qspec:eq:poisson} is independent of $a$.
Differentiate its remaining series and use
$(s-1)\Gamma(s-1)=\Gamma(s)$ to obtain
\eqref{qspec:eq:odd-poisson}. Normal convergence proves that this expression
is entire. The apparent singularity in \eqref{qspec:eq:odd-differentiation}
at $s=1$ is removable, not a missing residue.
\end{proof}

\subsection{The two local spectral germs}
For $\epsilon$ near zero define
\begin{align}
C(\epsilon;b)&=\frac{\sqrt\pi\,\Gamma(\epsilon)b^{-2\epsilon}}
                         {\Gamma(1/2+\epsilon)},\label{qspec:eq:Cgerm}\\
B_+(\epsilon;a,b)&=
\frac{4\pi^{1/2+\epsilon}b^{-\epsilon}}{\Gamma(1/2+\epsilon)}
\sum_{n\ge1}n^\epsilon\cos(2\pi na)K_\epsilon(2\pi bn),\label{qspec:eq:Bplusgerm}\\
B_-(\epsilon;a,b)&=
\frac{4\pi^{1+\epsilon}b^{1/2-\epsilon}}{\Gamma(1+\epsilon)}
\sum_{n\ge1}n^{1/2+\epsilon}\sin(2\pi na)
 K_{\epsilon-1/2}(2\pi bn).\label{qspec:eq:Bminusgerm}
\end{align}
Then
\begin{equation}\label{qspec:eq:germ-identities}
E(\tfrac12+\epsilon)=C(\epsilon;b)+B_+(\epsilon;a,b),
\qquad Y(1+\epsilon)=B_-(\epsilon;a,b).
\end{equation}
The only pole of these germs is the unit-residue pole of $C$ at zero.
Define their regular jets by
\begin{equation}\label{qspec:eq:jet-definitions}
C_m(b)=m![\epsilon^m]\bigl(C(\epsilon;b)-\epsilon^{-1}\bigr),
\qquad B_m^\pm(a,b)=\left.\partial_\epsilon^m B_\pm(\epsilon;a,b)\right|_{0}.
\end{equation}
The spectral shifts $1/2+\epsilon$ and $1+\epsilon$ differ deliberately:
the first-moment factor in $Y$ is what aligns both massless powers with
$x^{-1-2\epsilon}$.

% END sections/02-transform

% BEGIN sections/03-stieltjes
\section{Complete Stieltjes reflection through a positive mass}
\label{qspec:sec:stieltjes}

Define the following ordinary, absolutely convergent sums:
\begin{align}
D_m^+(a,b)=\sum_{n\ge0}\biggl\{&
\frac{(-\log Q_b(x_n))^m}{\sqrt{Q_b(x_n)}}
-\frac{(-2\log x_n)^m}{x_n}\notag\\[-2pt]
&+\frac{(-\log Q_b(y_n))^m}{\sqrt{Q_b(y_n)}}
-\frac{(-2\log y_n)^m}{y_n}\biggr\},\label{qspec:eq:Dplus}\\
D_m^-(a,b)=\sum_{n\ge0}\biggl\{&
\frac{x_n(-\log Q_b(x_n))^m}{Q_b(x_n)}
-\frac{(-2\log x_n)^m}{x_n}\notag\\[-2pt]
&-\frac{y_n(-\log Q_b(y_n))^m}{Q_b(y_n)}
+\frac{(-2\log y_n)^m}{y_n}\biggr\}.
\label{qspec:eq:Dminus}
\end{align}
The cancellation takes place separately at $x_n$ and $y_n$; no pairing of
two divergent sums is being used to define these differences.

\begin{theorem}[All-index, both-parity Stieltjes identity]
\label{qspec:thm:master}
For every $m\ge0$, $0<a<1$, and $b>0$,
\begin{align}
(-2)^mG_m^+(a)&=C_m(b)+B_m^+(a,b)-D_m^+(a,b),\label{qspec:eq:master-plus}\\
(-2)^mG_m^-(a)&=B_m^-(a,b)-D_m^-(a,b).\label{qspec:eq:master-minus}
\end{align}
In particular each individual Stieltjes constant is given by
\begin{equation}\label{qspec:eq:individual}
\gamma_m(a)=\frac{C_m+B_m^++B_m^--D_m^+-D_m^-}{2(-2)^m}.
\end{equation}
These identities may be differentiated any fixed number of times in $a$
and $b$ on compact subsets of \eqref{qspec:eq:domain}.
\end{theorem}
\begin{proof}
For $|\Rea\epsilon|<1/4$ put
\begin{align*}
D_+(\epsilon)&=\sum_{n\ge0}\left\{
Q_b(x_n)^{-1/2-\epsilon}-x_n^{-1-2\epsilon}
+Q_b(y_n)^{-1/2-\epsilon}-y_n^{-1-2\epsilon}\right\},\\
D_-(\epsilon)&=\sum_{n\ge0}\left\{
x_nQ_b(x_n)^{-1-\epsilon}-x_n^{-1-2\epsilon}
-y_nQ_b(y_n)^{-1-\epsilon}+y_n^{-1-2\epsilon}\right\}.
\end{align*}
Writing $Q_b(x)=x^2(1+b^2/x^2)$ shows that each local difference
is $O(x^{-3+2|\Rea\epsilon|})$, uniformly on compact sets of the
parameters and of $\epsilon$. Differentiation in $\epsilon$ introduces
at most fixed powers of $\log x$. The resulting series converge normally;
the same reasoning applies after finitely many $a,b$ derivatives.
Their $m$th derivatives at zero are exactly
\eqref{qspec:eq:Dplus} and~\eqref{qspec:eq:Dminus}.

In the initially convergent region, and hence as local meromorphic germs,
\begin{align}
E(\tfrac12+\epsilon)&=
\zeta(1+2\epsilon,a)+\zeta(1+2\epsilon,1-a)+D_+(\epsilon),\label{qspec:eq:even-subtraction}\\
Y(1+\epsilon)&=
\zeta(1+2\epsilon,a)-\zeta(1+2\epsilon,1-a)+D_-(\epsilon).
\label{qspec:eq:odd-subtraction}
\end{align}
Each individual Hurwitz term has residue $1/(2\epsilon)$, so the even
sum has residue $1/\epsilon$ and the odd difference has none. Equation
\eqref{qspec:eq:stieltjes-convention} gives the $m$th regular derivative
as $(-2)^mG_m^\pm$. Substitute \eqref{qspec:eq:germ-identities}, remove
the one explicit pole, and compare derivatives. This proves both
identities and their differentiated versions. Adding them proves
\eqref{qspec:eq:individual}.
\end{proof}

\subsection{All coefficients of the pole term}
The zero-mode coefficients have a particularly short finite description.
Set
\begin{equation}\label{qspec:eq:ell}
\ell_1=2\log(2/b),\qquad
\ell_r=(-1)^r(r-1)!(2-2^r)\zeta(r)\quad(r\ge2).
\end{equation}
\begin{proposition}\label{qspec:prop:pole-bell}
For every $m\ge0$,
\begin{equation}\label{qspec:eq:Cbell}
C_m(b)=\frac{\Bcal_{m+1}(\ell_1,\ldots,\ell_{m+1})}{m+1}.
\end{equation}
Writing $L=\log(2/b)$, the first three are
\begin{equation}\label{qspec:eq:Cexamples}
C_0=2L,\qquad C_1=2L^2-\zeta(2),\qquad
C_2=\frac83L^3-\frac{2\pi^2}{3}L+4\zeta(3).
\end{equation}
\end{proposition}
\begin{proof}
The function $\epsilon C(\epsilon;b)$ is analytic and equals one at zero.
Its logarithm is
\[
\log\Gamma(1+\epsilon)-\log\Gamma(1/2+\epsilon)
+\tfrac12\log\pi-2\epsilon\log b.
\]
The classical polygamma values at $1$ and $1/2$ \cite{gamma} give precisely
\eqref{qspec:eq:ell} as its successive derivatives at zero. Expand its
exponential by \eqref{qspec:eq:bell}, divide by $\epsilon$, remove the
pole, and multiply the coefficient of $\epsilon^m$ by $m!$. The resulting
factor is $m!/(m+1)!=1/(m+1)$, proving the proposition.
\end{proof}

\subsection{Explicit Fourier--Bessel spectral jets}
For $n\ge1$ set
\begin{align}
v_1(n)&=\gamma+\log(4\pi n/b),&
v_r&=(-1)^{r+1}(r-1)!(2^r-1)\zeta(r)\quad(r\ge2),\label{qspec:eq:vjets}\\
w_1(n)&=\gamma+\log(\pi n/b),&
w_r&=(-1)^{r+1}(r-1)!\zeta(r)\quad(r\ge2).\label{qspec:eq:wjets}
\end{align}
Define $K^{\langle r\rangle}_\nu(u)=\partial_\nu^rK_\nu(u)$.
Differentiating the holomorphic factors in
\eqref{qspec:eq:Bplusgerm}--\eqref{qspec:eq:Bminusgerm} gives
\begin{align}
B_m^+={}&4\sum_{n\ge1}\cos(2\pi na)
 \sum_{r=0}^{\lfloor m/2\rfloor}\binom{m}{2r}
 \Bcal_{m-2r}(v(n))K^{\langle2r\rangle}_0(2\pi bn),\label{qspec:eq:Bplusjets}\\
B_m^-={}&4\pi\sqrt b\sum_{n\ge1}\sqrt n\sin(2\pi na)
 \sum_{r=0}^m\binom mr\Bcal_{m-r}(w(n))
 K^{\langle r\rangle}_{-1/2}(2\pi bn).
\label{qspec:eq:Bminusjets}
\end{align}
Here $\Bcal_j(v(n))$ means that only the first $j$ entries are used.
Equation~\eqref{qspec:eq:K-cosh} shows that $K_\nu=K_{-\nu}$, so all
odd order derivatives at zero vanish. No corresponding deletion is valid
at $\nu=-1/2$. This explains the different inner sums.

For example,
\[
B_0^+=4\sum_{n\ge1}\cos(2\pi na)K_0(2\pi bn),\qquad
B_1^+=4\sum_{n\ge1}\cos(2\pi na)
\bigl(\gamma+\log(4\pi n/b)\bigr)K_0(2\pi bn).
\]
Thus \eqref{qspec:eq:master-plus} at $m=1$ is the explicit identity
\begin{align}
-2G_1^+(a)={}&2\log^2(2/b)-\zeta(2)-D_1^+(a,b)\notag\\
&+4\sum_{n\ge1}\cos(2\pi na)
\bigl(\gamma+\log(4\pi n/b)\bigr)K_0(2\pi bn).
\label{qspec:eq:first-stieltjes-example}
\end{align}
It involves no Bessel order derivative. Higher $m$ introduce them only
at the even orders displayed in \eqref{qspec:eq:Bplusjets}.

\subsection{Damped polylogarithm integral formulas}
Put
\begin{equation}\label{qspec:eq:zt}
z_t^\pm=\exp(-2\pi b\cosh t\pm2\pi ia),\qquad
\mathcal C_\sigma(t)=\frac{\Li_\sigma(z_t^+)+\Li_\sigma(z_t^-)}2,
\quad
\mathcal S_\sigma(t)=\frac{\Li_\sigma(z_t^+)-\Li_\sigma(z_t^-)}{2i}.
\end{equation}
For real $\sigma$, these are respectively the real and imaginary parts of
$\Li_\sigma(z_t^+)$. For complex $\sigma$, the symmetric formulas in
\eqref{qspec:eq:zt}, not a real-part operation, preserve holomorphy.

\begin{theorem}[Polylogarithm order integrals for both reflected channels]
\label{qspec:thm:polylog-integrals}
The entire germs $B_\pm$ satisfy
\begin{align}
B_+(\epsilon;a,b)&=
\frac{4\pi^{1/2+\epsilon}b^{-\epsilon}}{\Gamma(1/2+\epsilon)}
\int_0^\infty\cosh(\epsilon t)\mathcal C_{-\epsilon}(t)\dd t,
\label{qspec:eq:polylogplus}\\
B_-(\epsilon;a,b)&=
\frac{4\pi^{1+\epsilon}b^{1/2-\epsilon}}{\Gamma(1+\epsilon)}
\int_0^\infty\cosh((\epsilon-\tfrac12)t)
\mathcal S_{-\epsilon-1/2}(t)\dd t.
\label{qspec:eq:polylogminus}
\end{align}
Every spectral derivative in Theorem~\ref{qspec:thm:master} can therefore
be represented by differentiating an ordinary absolutely convergent
polylogarithm integral.
\end{theorem}
\begin{proof}
Insert \eqref{qspec:eq:K-cosh} in the Bessel series. The polylogarithm
series \cite{polylog} gives
\[
\mathcal C_{-\epsilon}(t)=\sum_{n\ge1}n^\epsilon
 e^{-2\pi bn\cosh t}\cos(2\pi na),
\]
with the analogous sine formula containing $n^{\epsilon+1/2}$.
For $\epsilon$ in a compact set, the terms and their derivatives have
majorants of the form
\[
C n^A(1+\log n)^M(1+t)^M e^{At}e^{-2\pi b_0n\cosh t},
\qquad b\ge b_0>0.
\]
Summing in $n$ and integrating in $t$ gives a finite quantity. The decay
in $t$ is stronger than any exponential. Fubini's theorem and holomorphic
parameter differentiation therefore justify the interchange and all
claimed spectral derivatives.
\end{proof}

\subsection{Zeroth-order reflection and an exact mass-transfer law}
Since $K_{1/2}(u)=\sqrt{\pi/(2u)}e^{-u}$,
\begin{equation}\label{qspec:eq:odd-zero}
B_0^-(a,b)=2\pi\sum_{n\ge1}e^{-2\pi bn}\sin(2\pi na)
=\frac{\pi\sin(2\pi a)}{\cosh(2\pi b)-\cos(2\pi a)}.
\end{equation}
The odd master identity recovers the classical digamma reflection law:
\begin{equation}\label{qspec:eq:zero-reflection}
\pi\cot(\pi a)=
\frac{\pi\sin(2\pi a)}{\cosh(2\pi b)-\cos(2\pi a)}-D_0^-(a,b).
\end{equation}
At $a=1/2$ every odd term is exactly zero, a structural check rather
than a near-zero numerical test.

Because the left sides of the master identities are independent of $b$,
we also obtain, for arbitrary $b,c>0$,
\begin{align}
D_m^+(a,b)-D_m^+(a,c)&=C_m(b)-C_m(c)+B_m^+(a,b)-B_m^+(a,c),\label{qspec:eq:transferplus}\\
D_m^-(a,b)-D_m^-(a,c)&=B_m^-(a,b)-B_m^-(a,c).
\label{qspec:eq:transferminus}
\end{align}
These are exact identities between two deformations with the Stieltjes
constant eliminated. They generate harmonic/polylogarithm identities
without retaining an unevaluated Stieltjes symbol.

% END sections/03-stieltjes

% BEGIN sections/04-harmonic-primitives
\section{Harmonic coefficients, global tails, and parameter primitives}
\label{qspec:sec:harmonic}

\subsection{The ordinary and half-shifted harmonic dictionaries}
For $\alpha>0$ define
\begin{equation}\label{qspec:eq:shifted-harmonics}
h_{k,r}^{(\alpha)}=\sum_{j=0}^{k-1}(j+\alpha)^{-r},\qquad
 d_{k,r}^{(\alpha)}=(-1)^{r-1}(r-1)!h_{k,r}^{(\alpha)}.
\end{equation}
The two cases needed in the present problem are
\begin{equation}\label{qspec:eq:harmonic-dictionaries}
h_{k,r}^{(1)}=H_k^{(r)},\qquad
h_{k,r}^{(1/2)}=2^r H_{2k}^{(r)}-H_k^{(r)}.
\end{equation}
These follow just by separating even and odd integers. Logarithmic
differentiation of a finite Pochhammer product gives the exact finite law
\begin{equation}\label{qspec:eq:pochhammer-jet}
\left.\partial_\epsilon^j(\alpha+\epsilon)_k\right|_0
=(\alpha)_k\Bcal_j\bigl(d_{k,1}^{(\alpha)},\ldots,d_{k,j}^{(\alpha)}\bigr).
\end{equation}
In particular, the first four factors after dividing by $(\alpha)_k$ are
\[
1,\qquad h_{k,1},\qquad h_{k,1}^2-h_{k,2},\qquad
h_{k,1}^3-3h_{k,1}h_{k,2}+2h_{k,3}.
\]
The signs are those of a rising Pochhammer product, not its reciprocal.

\begin{theorem}[All-index harmonic Hurwitz counterterms]
\label{qspec:thm:harmonic}
Let $\eta\in\{+1,-1\}$ and put $\alpha_{+1}=1/2$, $\alpha_{-1}=1$.
If $0<b<\min(a,1-a)$, then
\begin{align}
D_m^\eta(a,b)={}&\sum_{k\ge1}
 \frac{(-1)^k(\alpha_\eta)_k b^{2k}}{k!}
 \sum_{j=0}^m\binom mj 2^{m-j}\Bcal_j(d_k^{(\alpha_\eta)})\notag\\
&\quad\times\left\{\zeta^{(m-j)}(2k+1,a)
       +\eta\zeta^{(m-j)}(2k+1,1-a)\right\}.
\label{qspec:eq:harmonic-counterterms}
\end{align}
This is an absolutely and locally uniformly convergent identity, including
all fixed-order parameter derivatives. For the even channel,
$(1/2)_k/k!=\binom{2k}{k}/4^k$; for the odd channel $(1)_k/k!=1$.
\end{theorem}
\begin{proof}
The common local factor in the counterterm germs is
\[
x^{-1-2\epsilon}\left((1+b^2/x^2)^{-\alpha_\eta-\epsilon}-1\right)
=\sum_{k\ge1}\frac{(-1)^k(\alpha_\eta+\epsilon)_k b^{2k}}{k!}
 x^{-2k-1-2\epsilon}.
\]
The binomial expansion converges uniformly with a geometric ratio strictly
less than one for $x=x_n,y_n$ and the stated parameter range. Summing in
$n$ gives the two Hurwitz terms. Local uniform convergence also follows
directly by bounding each Hurwitz tail with its leading geometric scale;
polynomial factors in $k$ and $\log k$, introduced by differentiation,
are harmless on a smaller compact set. The product rule, the chain factor
$2^{m-j}$, and \eqref{qspec:eq:pochhammer-jet} yield
\eqref{qspec:eq:harmonic-counterterms}.
\end{proof}

\subsection{Finite-head continuation for arbitrary positive mass}
The small-mass hypothesis in Theorem~\ref{qspec:thm:harmonic} is essential
for that particular unshifted power series. It is not a restriction on
the Stieltjes identities.

\begin{theorem}[An ordinary convergent finite-head formula]
\label{qspec:thm:head}
Choose any integer $M\ge0$ such that
\begin{equation}\label{qspec:eq:head-condition}
b<\min(M+a,M+1-a).
\end{equation}
Let $D_{m,<M}^\eta$ be the sum of the first $M$ summands in
\eqref{qspec:eq:Dplus} or~\eqref{qspec:eq:Dminus}. Then
\begin{align}
D_m^\eta(a,b)={}&D_{m,<M}^\eta(a,b)
+\sum_{k\ge1}\frac{(-1)^k(\alpha_\eta)_k b^{2k}}{k!}
\sum_{j=0}^m\binom mj2^{m-j}\Bcal_j(d_k^{(\alpha_\eta)})\notag\\
&\quad\times\left\{\zeta^{(m-j)}(2k+1,M+a)
+\eta\zeta^{(m-j)}(2k+1,M+1-a)\right\}.
\label{qspec:eq:finite-head}
\end{align}
For every $a,b$ in \eqref{qspec:eq:domain}, such an $M$ exists, and the
formula converges normally in a neighborhood of those parameters.
\end{theorem}
\begin{proof}
Remove the finite head before performing the binomial expansion. In the
remaining sum the smallest positive arguments are $M+a$ and $M+1-a$,
so \eqref{qspec:eq:head-condition} is exactly the required geometric
condition. The proof of Theorem~\ref{qspec:thm:harmonic} then applies to
each tail independently. Increasing $M$ only transfers exact summands
between the head and the tails; it changes neither the function nor any
of its jets.
\end{proof}

There is an important bookkeeping point: the second tail argument is
$M+1-a$, \emph{not} $1-(M+a)$. Reflection describes the original two
one-sided spectra; it does not reverse the direction of the index shift.
For an implementation one can choose $M$ with
$\min(M+a,M+1-a)\ge2b$. Then both power-series ratios are at most $1/4$.
This choice is used by the supplied numerical code; the mathematical
identity does not depend on this particular convenience.

For example, the even zeroth counterterm is
\[
D_0^+(a,b)=\sum_{k\ge1}(-1)^k\frac{\binom{2k}{k}}{4^k}b^{2k}
\bigl(\zeta(2k+1,a)+\zeta(2k+1,1-a)\bigr)
\]
inside the small-mass disk. Its first spectral derivative inserts
\[
\bigl(2H_{2k}-H_k\bigr)
\bigl(\zeta(2k+1,a)+\zeta(2k+1,1-a)\bigr)
+2\bigl(\zeta'(2k+1,a)+\zeta'(2k+1,1-a)\bigr).
\]
Replacing the two Hurwitz arguments by their shifted tail arguments and
adding the finite head makes both identities valid for every $b>0$.

\subsection{Every fixed parameter derivative}
The classical identity
$\partial_a^p\zeta(s,a)=(-1)^p(s)_p\zeta(s+p,a)$ \cite{hurwitz}
produces a second finite harmonic dictionary.

\begin{proposition}\label{qspec:prop:parameter-derivatives}
For $p\ge1$ and $m\ge0$, put
$d_{p,r}=(-1)^{r-1}(r-1)!H_p^{(r)}$. Then
\begin{equation}\label{qspec:eq:stieltjes-derivative}
\gamma_m^{(p)}(a)=(-1)^{m+p}p!
\sum_{j=0}^m\binom mj\Bcal_j(d_p)
\zeta^{(m-j)}(p+1,a).
\end{equation}
\end{proposition}
\begin{proof}
Differentiate the Laurent expansion
\eqref{qspec:eq:stieltjes-convention} in $a$. Its pole is independent of
$a$ and disappears. Expand $(1+u)_p$ by
\eqref{qspec:eq:pochhammer-jet} with $\alpha=1$, multiply by
$\zeta(p+1+u,a)$, and compare the coefficient of $u^m$.
\end{proof}
For $m=0$ this gives $\gamma_0^{(p)}(a)=(-1)^pp!\zeta(p+1,a)$,
consistent with $\gamma_0=-\psi$. Parameter differentiations of the
master identity therefore produce exact harmonic identities involving
positive-order polygamma and spectral Hurwitz derivatives, without
introducing a new class of constants.

\subsection{First and repeated anchored parameter primitives}
A convenient first primitive is
\begin{equation}\label{qspec:eq:Pprimitive}
P_m(a)=\frac{(-1)^{m+1}}{m+1}\zeta^{(m+1)}(0,a),
\qquad P_m'(a)=\gamma_m(a).
\end{equation}
Indeed, differentiating $\partial_a\zeta(u,a)=-u\zeta(1+u,a)$
$m+1$ times in $u$ gives this identity. At $m=0$ it is
$P_0=-\log\Gamma+\tfrac12\log(2\pi)$, the usual Lerch relation.
Consequently the primitives of the two reflection channels anchored at
$a=1/2$ are
\begin{align}
\int_{1/2}^a G_m^+(t)\dd t&=P_m(a)-P_m(1-a),\label{qspec:eq:plus-primitive}\\
\int_{1/2}^a G_m^-(t)\dd t&=P_m(a)+P_m(1-a)-2P_m(1/2).
\label{qspec:eq:minus-primitive}
\end{align}
All integrals of the counterterm series over such compact intervals can
be performed termwise by normal convergence.

For completeness, every higher anchored primitive also has a finite
Hurwitz-jet description. Let
\[
(T_{r-1,a_0}f)(a)=\sum_{j=0}^{r-1}\frac{(a-a_0)^j}{j!}f^{(j)}(a_0),
\qquad r\ge1,\quad a,a_0>0,
\]
and define, initially for $u\ne0$ near zero,
\begin{equation}\label{qspec:eq:anchored-germ}
\mathcal I_r(u;a,a_0)=
\frac{(1-T_{r-1,a_0})\zeta(1+u-r,a)}{(-u)_r}
-\frac{(a-a_0)^r}{r!u}.
\end{equation}
Here the Taylor operator acts on the parameter $a$, not on $u$.

\begin{proposition}[All anchored primitive orders]
\label{qspec:prop:all-primitives}
The apparent pole in \eqref{qspec:eq:anchored-germ} is removable. For
$m\ge0$,
\begin{equation}\label{qspec:eq:all-primitives}
\frac1{(r-1)!}\int_{a_0}^a(a-t)^{r-1}\gamma_m(t)\dd t
=(-1)^m\left.\partial_u^m\mathcal I_r(u;a,a_0)\right|_{u=0}.
\end{equation}
The oriented integral convention also covers $a<a_0$.
\end{proposition}
\begin{proof}
The numerator at $u=0$ is the Taylor remainder of
$\zeta(1-r,a)=-B_r(a)/r$, a degree-$r$ polynomial with leading
coefficient $-1/r$. Its value is $-(a-a_0)^r/r$.
As $(-u)_r=-(r-1)!u+O(u^2)$, the residue of the quotient is exactly
the term explicitly subtracted in \eqref{qspec:eq:anchored-germ}.
The remainder is holomorphic in $u$.

Further,
\[
\partial_a^r\zeta(1+u-r,a)=(-u)_r\zeta(1+u,a).
\]
Taking $r$ parameter derivatives of \eqref{qspec:eq:anchored-germ}
therefore yields $\zeta(1+u,a)-1/u$. All its first $r$ anchored
initial values, through derivative $r-1$, vanish. Uniqueness of the
anchored repeated primitive and the Laurent expansion prove
\eqref{qspec:eq:all-primitives}.
\end{proof}
This is a classical Hurwitz-jet integration mechanism written in a
pole-cancelled form. Coefficient extraction requires only finitely many
Hurwitz spectral derivatives and Pochhammer coefficients at each fixed
$m,r$.

\subsection{An explicit log-Gamma ratio identity}
A zeroth-order primitive of the even master produces a particularly
simple ordinary sum. Define
\begin{equation}\label{qspec:eq:A}
A(a,b)=\sum_{n\ge0}\left\{
\arsinh\frac{n+a}{b}-\arsinh\frac{n+1-a}{b}
-\log\frac{n+a}{n+1-a}\right\}.
\end{equation}
The local primitive
$\arsinh(x/b)-\log x=\log(2/b)+O(x^{-2})$ proves convergence of this
paired expression and its needed parameter derivatives.

\begin{corollary}\label{qspec:cor:gamma-ratio}
For $0<a<1$, $b>0$,
\begin{align}
\log\frac{\Gamma(1-a)}{\Gamma(a)}={}&
2(a-\tfrac12)\log(2/b)-A(a,b)\notag\\
&+\frac2\pi\sum_{n\ge1}\frac{\sin(2\pi na)}nK_0(2\pi bn).
\label{qspec:eq:gamma-ratio}
\end{align}
If $M$ satisfies \eqref{qspec:eq:head-condition}, the tail of $A$ is
\begin{equation}\label{qspec:eq:A-tail}
A_{\ge M}(a,b)=\sum_{k\ge1}
\frac{(-1)^{k+1}(1/2)_k b^{2k}}{2k\,k!}
\bigl(\zeta(2k,M+a)-\zeta(2k,M+1-a)\bigr).
\end{equation}
\end{corollary}
\begin{proof}
Integrate \eqref{qspec:eq:master-plus} at $m=0$ from $1/2$ to $a$.
The left side is \eqref{qspec:eq:plus-primitive} at $m=0$.
The pole term integrates to $2(a-1/2)\log(2/b)$; the cosine terms
integrate to the sine series. Differentiating \eqref{qspec:eq:A} in $a$
gives $D_0^+(a,b)$, and $A(1/2,b)=0$. This proves
\eqref{qspec:eq:gamma-ratio}. Finally integrate the binomial expansion
of $Q_b(x)^{-1/2}-x^{-1}$ in $x$ and subtract its two tail values.
The coefficients are precisely those in \eqref{qspec:eq:A-tail}.
\end{proof}

% END sections/04-harmonic-primitives

% BEGIN sections/05-anomaly
\section{Conjugate factorization and an exact polynomial parity criterion}
\label{qspec:sec:anomaly}

Although $\log(x^2+b^2)=\log(x+ib)+\log(x-ib)$ for $x>0$, one cannot
split arbitrary zeta-regularized weighted sums of these logarithms without
checking the continuation. This section computes the correction rather
than assuming it vanishes.

For $a,b>0$ and a polynomial $W(x)=\sum_{d=0}^D w_dx^d\in\C[x]$, define
\begin{align}
Z_W(s;a,b)&=\sum_{n\ge0}W(n+a)Q_b(n+a)^{-s},\label{qspec:eq:weighted-quadratic}\\
F_{W,\pm}(s;a,b)&=\sum_{n\ge0}W(n+a)(n+a\pm ib)^{-s}.
\label{qspec:eq:weighted-linear}
\end{align}
Their initial absolute-convergence conditions are respectively
$\Rea s>(D+1)/2$ and $\Rea s>D+1$. We always mean their meromorphic
continuations when evaluating at zero. The sign convention for the
factorization defect is
\begin{equation}\label{qspec:eq:delta-definition}
\delta_W(a,b)=F_{W,+}'(0;a,b)+F_{W,-}'(0;a,b)-Z_W'(0;a,b).
\end{equation}
It is a convention for spectral derivatives; changing to logarithms of
regularized determinants changes the associated sign convention.

\begin{theorem}[Harmonic residue formula and universal vanishing criterion]
\label{qspec:thm:anomaly}
All three continued functions in
\eqref{qspec:eq:weighted-quadratic}--\eqref{qspec:eq:weighted-linear} are
holomorphic at zero. The defect is independent of $a$ and equals
\begin{align}
\delta_W(a,b)
&=\frac14[x^{-1}]W(x)
\left(\log(1+ib/x)-\log(1-ib/x)\right)^2\label{qspec:eq:anomaly-residue}\\
&=\sum_{k\ge1}w_{2k-1}\frac{(-1)^k b^{2k}}{k}
\left(H_{2k}-\frac12H_k\right).
\label{qspec:eq:anomaly-harmonic}
\end{align}
The last sum is finite. For a fixed polynomial $W$, the identity
$\delta_W(a,b)=0$ for every $b>0$ holds if and only if $W$ is even.
\end{theorem}
\begin{proof}
We first take $b<a$. Put $\ell_\pm(x)=\log(1\pm ib/x)$ and
$L=\ell_++\ell_-$. In a common initial half-plane, factor out $x^{-2s}$
and compare $Z_W(s)$ with
$\tfrac12(F_{W,+}(2s)+F_{W,-}(2s))$. The remaining difference is
\begin{equation}\label{qspec:eq:local-defect}
e^{-sL}-\tfrac12\bigl(e^{-2s\ell_+}+e^{-2s\ell_-}\bigr)
=-\frac{s^2}{2}(\ell_+-\ell_-)^2+O(s^3).
\end{equation}
Both the constant and first powers of $s$ cancel identically as formal
power series in $x^{-1}$.

To justify the continuation and the coefficient comparison, expand each
factor in powers of $x^{-1}$ and sum $x=n+a$. A monomial $w_dx^d$
paired with the coefficient of $x^{-j}$ contributes
$w_d\zeta(2s+j-d,a)$. For $s$ near zero, only $j=d+1$ has a pole,
and it is
\begin{equation}\label{qspec:eq:hurwitz-local-pole}
\zeta(1+2s,a)=\frac1{2s}+\gamma_0(a)+O(s).
\end{equation}
Every original positive-degree coefficient in the exponential expansion
vanishes at least to order one at $s=0$, so this possible pole is
cancelled. The degree-zero term is a finite combination of $\zeta(-d,a)$,
which is regular. The remaining infinite coefficient sum converges
normally near zero because $b/a<1$. This proves holomorphy of the
original functions at zero, not just holomorphy of their difference.

In the difference \eqref{qspec:eq:local-defect}, all regular Hurwitz
terms start at $s^2$ and have zero first derivative at zero. Multiplying
the pole in \eqref{qspec:eq:hurwitz-local-pole} by the quadratic term
in \eqref{qspec:eq:local-defect} gives
\[
Z_W'(0)-F_{W,+}'(0)-F_{W,-}'(0)
=-\tfrac14[x^{-1}]W(x)(\ell_+-\ell_-)^2.
\]
The finite part $\gamma_0(a)$ contributes only from order $s^2$ onward;
this explains both the sign and the independence from $a$.

For arbitrary $b>0$, remove a finite head with $M+a>b$ before expanding.
At each of the finitely many head terms, the first derivatives cancel
because the ordinary logarithms add with the stated branches. The tail
calculation is identical, and the residue in
\eqref{qspec:eq:hurwitz-local-pole} is still one in its spectral variable.
This proves \eqref{qspec:eq:anomaly-residue} globally.

For the coefficient formula, write
\[
\ell_+-\ell_-=2\sum_{r\ge0}\frac{(ib)^{2r+1}}{(2r+1)x^{2r+1}}.
\]
The coefficient of $x^{-2k}$ in one quarter of its square is
\[
(-1)^kb^{2k}\sum_{r=0}^{k-1}
\frac1{(2r+1)(2k-2r-1)}
=\frac{(-1)^kb^{2k}}k\sum_{r=0}^{k-1}\frac1{2r+1}.
\]
The last sum is $H_{2k}-H_k/2$. It is nonzero for every $k\ge1$.
Only odd monomials $x^{2k-1}$ can contribute to $[x^{-1}]$.
Finally the distinct powers $b^{2k}$ are linearly independent as
polynomials in $b$, proving the universal vanishing criterion.
\end{proof}

\begin{remark}[A fixed mass is different from all masses]
At a single prescribed value of $b$, different odd terms can cancel.
For example $W(x)=x^3+\tfrac23b_0^2x$ has zero defect at $b=b_0$
but is not even. Thus the theorem characterizes vanishing identically
in the mass, not an erroneous pointwise equivalence.
\end{remark}

The first defect values, with the convention
\eqref{qspec:eq:delta-definition}, are
\begin{equation}\label{qspec:eq:anomaly-examples}
\delta_1=0,\qquad \delta_x=-b^2,\qquad
\delta_{x^2}=0,\qquad \delta_{x^3}=\frac23b^4,\qquad
\delta_{x^5}=-\frac{23}{45}b^6.
\end{equation}
They are useful negative controls: deleting the anomaly from an odd-weight
calculation gives an explicitly false formula, not merely an uncertified
one.

\subsection{Finite corrected formulas for any polynomial weight}
Expand $W$ about either conjugate root by writing
\begin{equation}\label{qspec:eq:shifted-weight-coeffs}
W(y\mp ib)=\sum_{r=0}^D c_r^\pm(b)y^r.
\end{equation}
The linear-factor series are finite sums of shifted Hurwitz functions,
so Theorem~\ref{qspec:thm:anomaly} immediately gives
\begin{equation}\label{qspec:eq:general-weighted-hurwitz}
Z_W'(0;a,b)=
\sum_{r=0}^D\left(c_r^+(b)\zeta'(-r,a+ib)
+c_r^-(b)\zeta'(-r,a-ib)\right)-\delta_W(a,b).
\end{equation}
For real coefficients, the two finite sums are conjugates. Thus every
polynomially weighted quadratic logarithmic sum in this class has a
finite Hurwitz-derivative evaluation, with the exact harmonic correction
already displayed. This is a concrete specialization of the classical
local-anomaly principle, not an assumption of multiplicativity.

\subsection{Anomaly-free factorization at every negative integer}
Let
\begin{equation}\label{qspec:eq:one-sided}
Z_a(s;b)=\sum_{n\ge0}Q_b(n+a)^{-s}
\end{equation}
initially for $\Rea s>1/2$. The weight $W(x)=Q_b(x)^m$ is even for every
integer $m\ge0$, so the defect vanishes.

\begin{corollary}\label{qspec:cor:negative-hurwitz}
For $a,b>0$ and $m\ge0$,
\begin{equation}\label{qspec:eq:negative-hurwitz}
Z_a'(-m;b)=2\Rea\sum_{j=0}^m
\binom mj(-2ib)^{m-j}\zeta'(-m-j,a+ib).
\end{equation}
\end{corollary}
\begin{proof}
For this weight $Z_W(s)=Z_a(s-m;b)$. With $y=x+ib$,
\[
Q_b(x)^m=y^m(y-2ib)^m
=\sum_{j=0}^m\binom mj(-2ib)^{m-j}y^{m+j}.
\]
Insert this finite expansion into
\eqref{qspec:eq:general-weighted-hurwitz}; the defect is zero and the
second conjugate factor gives the real-part formula.
\end{proof}
No claim is made that an arbitrary higher spectral derivative factorizes
in this way. The proof uses the first-derivative cancellation and the
special even polynomial weight. At higher spectral orders, logarithmic
cross terms require their own calculation.

% END sections/05-anomaly

% BEGIN sections/06-polylogarithms
\section{Finite polylogarithmic reductions and mass antiderivatives}
\label{qspec:sec:finite}

\subsection{All nonpositive integer spectral points}
Put
\begin{equation}\label{qspec:eq:Lm}
L_m(a,b)=\left.\partial_s E(s;a,b)\right|_{s=-m},\qquad
z=\exp(-2\pi b+2\pi ia),\qquad m\ge0.
\end{equation}
The half-integer Bessel expression is the finite polynomial
\begin{equation}\label{qspec:eq:half-integer-K}
K_{m+1/2}(u)=\sqrt{\frac\pi{2u}}e^{-u}
\sum_{j=0}^m\frac{(m+j)!}{j!(m-j)!(2u)^j}.
\end{equation}
Its $m=0$ case and the three-term order recurrence prove it by induction;
it is also the classical half-integer formula \cite{halfK}.

\begin{theorem}[Finite Hurwitz--polylogarithm identities at every negative point]
\label{qspec:thm:finite}
For $0<a<1$, $b>0$, and every $m\ge0$,
\begin{align}
L_m(a,b)={}&2\Rea\sum_{j=0}^m\binom mj(-2ib)^{m-j}
\bigl\{\zeta'(-m-j,a+ib)+\zeta'(-m-j,1-a+ib)\bigr\}\label{qspec:eq:L-Hurwitz}\\
={}&-\frac{2^{2m+1}(m!)^2\pi}{(2m+1)!}b^{2m+1}\notag\\
&+2(-1)^mm!\sum_{j=0}^m
\frac{(m+j)!b^{m-j}}{j!(m-j)!4^j\pi^{m+j}}
\Rea\Li_{m+j+1}(z).
\label{qspec:eq:L-polylog}
\end{align}
In particular, the equality of the two right sides is a finite special-function
identity; $L_m$ is not needed as an additional unevaluated coordinate.
\end{theorem}
\begin{proof}
The reflected decomposition of $E$ is
$E(s;a,b)=Z_a(s;b)+Z_{1-a}(s;b)$. Apply
Corollary~\ref{qspec:cor:negative-hurwitz} to the two terms to obtain
\eqref{qspec:eq:L-Hurwitz}.

For the independent evaluation, write the Poisson expression as
$E(s)=H(s)/\Gamma(s)$, where $H$ is holomorphic at $s=-m$. The zero
of the reciprocal Gamma function is simple and
\[
(1/\Gamma)'(-m)=(-1)^mm!.
\]
Consequently $E'(-m)=(-1)^mm!H(-m)$: derivatives of the Bessel order
and of the other factors in $H$ do not appear in this first jet. The
zero Fourier mode becomes
\[
(-1)^mm!\sqrt\pi\Gamma(-m-\tfrac12)b^{2m+1}
=-\frac{2^{2m+1}(m!)^2\pi}{(2m+1)!}b^{2m+1}.
\]
For the nonzero modes use $K_{-m-1/2}=K_{m+1/2}$ and
\eqref{qspec:eq:half-integer-K}. The square-root prefactor reduces to
$1/(2\sqrt{bn})$. The remaining $n$ powers are
$n^{-m-j-1}$, so their cosine sums are exactly
$\Rea\Li_{m+j+1}(z)$. Collecting the powers of $b$, $4$, and $\pi$
gives \eqref{qspec:eq:L-polylog}.
\end{proof}

The first four values, writing $\Rea\Li_j=\Rea\Li_j(z)$, are
\begin{align}
L_0={}&-\log\bigl(2(\cosh(2\pi b)-\cos(2\pi a))\bigr),\label{qspec:eq:L0}\\
L_1={}&-\frac{4\pi b^3}{3}-\frac{2b}{\pi}\Rea\Li_2
                       -\frac1{\pi^2}\Rea\Li_3,\label{qspec:eq:L1}\\
L_2={}&-\frac{16\pi b^5}{15}+\frac{4b^2}{\pi^2}\Rea\Li_3
              +\frac{6b}{\pi^3}\Rea\Li_4+\frac3{\pi^4}\Rea\Li_5,\label{qspec:eq:L2}\\
L_3={}&-\frac{32\pi b^7}{35}-\frac{12b^3}{\pi^3}\Rea\Li_4
              -\frac{36b^2}{\pi^4}\Rea\Li_5\notag\\
     &\hspace{19mm}-\frac{45b}{\pi^5}\Rea\Li_6
              -\frac{45}{2\pi^6}\Rea\Li_7.
\label{qspec:eq:L3}
\end{align}
For $m=0$, $\Li_1(z)=-\log(1-z)$ reduces the finite formula to
\eqref{qspec:eq:L0}. For $m=1$, the Hurwitz side is exactly
\eqref{qspec:eq:opening-example}.

\subsection{An exact Appell-type ladder in the squared mass}
Here the term ``Appell-type'' means the displayed derivative relation,
not an assertion that these functions are polynomials.

\begin{theorem}[Mass differentiation and every anchored mass primitive]
\label{qspec:thm:mass-ladder}
With $u=b^2$,
\begin{equation}\label{qspec:eq:mass-derivative}
\partial_u L_m(a,\sqrt u)=mL_{m-1}(a,\sqrt u)\quad(m\ge1),
\qquad
\partial_u L_0(a,\sqrt u)=-E(1;a,\sqrt u).
\end{equation}
The continuous zero-mass values are
\begin{align}
L_0(a,0)&=-\log(4\sin^2\pi a),\label{qspec:eq:massless0}\\
L_m(a,0)&=\frac{2(-1)^m(2m)!}{(2\pi)^{2m}}
\Rea\Li_{2m+1}(e^{2\pi ia})\quad(m\ge1).
\label{qspec:eq:masslessm}
\end{align}
For every $m\ge1$,
\begin{align}
L_m(a,b)={}&\sum_{j=0}^{m-1}\binom mj b^{2j}L_{m-j}(a,0)\notag\\
&+m\int_0^{b^2}(b^2-u)^{m-1}L_0(a,\sqrt u)\dd u.
\label{qspec:eq:mass-primitives}
\end{align}
Thus the finite polylogarithm formulas evaluate every integral on the
right side exactly.
\end{theorem}
\begin{proof}
Initially in an absolute-convergence half-plane,
$\partial_u E(s)=-sE(s+1)$. Continue meromorphically and differentiate
in $s$:
\[
\partial_u E'(s)=-E(s+1)-sE'(s+1).
\]
At $s=-m$, $m\ge1$, the first term is zero because
$E(1-m)=0$. This proves the first relation. At $s=0$ it gives the
second relation.

For fixed $a\in(0,1)$ there is a nonzero gap between the lattice and
zero. The binomial continuation, after finitely many Hurwitz
subtractions if needed, is analytic in $u$ near zero. Alternatively,
take the limit in the explicit finite formula: for $m\ge1$ only its
$j=m$ term survives, while the absolutely convergent polylogarithms
have continuous radial limits. This proves
\eqref{qspec:eq:masslessm}; \eqref{qspec:eq:massless0} follows directly
from \eqref{qspec:eq:L0}. Repeated integration of
\eqref{qspec:eq:mass-derivative}, with these initial values, gives
\eqref{qspec:eq:mass-primitives}. The coefficient of its integral is
$m!/(m-1)!=m$.
\end{proof}

For example, a non-elementary logarithmic integral is evaluated by
\begin{align}
&-\int_0^{b^2}\log\bigl(2(\cosh(2\pi\sqrt u)-\cos(2\pi a))\bigr)\dd u\notag\\
&\qquad=-\frac{4\pi b^3}{3}-\frac{2b}{\pi}\Rea\Li_2(z)
-\frac1{\pi^2}\Rea\Li_3(z)
+\frac1{\pi^2}\Rea\Li_3(e^{2\pi ia}).
\label{qspec:eq:log-mass-integral}
\end{align}
The higher-weight analogues follow by inserting $m=2,3,\ldots$ in
\eqref{qspec:eq:mass-primitives}. Their constants are fixed by the
zero-mass boundary values, not left as undetermined integration terms.

The same ladder provides an independent algebraic check of the finite
coefficients. Since $\partial_b\Li_j(z)=-2\pi\Li_{j-1}(z)$,
substitution of \eqref{qspec:eq:L-polylog} into
$\partial_bL_m=2bmL_{m-1}$ leaves a finite identity of rational
coefficients and powers of $b,\pi$. The delivered exact suite checks
these cancellations through $m=10$; the theorem itself holds for all
$m$ by the analytic proof.

\subsection{A differential equation in the shift and mass}
The same finite family is closed under a second differential relation.
\begin{proposition}\label{qspec:prop:PDE}
For every $m\ge0$,
\begin{equation}\label{qspec:eq:PDE}
\left(\partial_a^2+\partial_b^2-\frac{2m}{b}\partial_b\right)L_m(a,b)=0.
\end{equation}
Equivalently, with $u=b^2$,
\[
\bigl(\partial_a^2+4u\partial_u^2+(2-4m)\partial_u\bigr)
L_m(a,\sqrt u)=0.
\]
For $m\ge2$ this becomes the finite adjacent-weight identity
\begin{equation}\label{qspec:eq:PDE-adjacent}
\partial_a^2L_m=2m(2m-1)L_{m-1}-4m(m-1)b^2L_{m-2}.
\end{equation}
\end{proposition}
\begin{proof}
Direct differentiation of one radial summand gives
\[
(\partial_a^2+\partial_b^2)E(s)=4s^2E(s+1),
\qquad \partial_bE(s)=-2sbE(s+1).
\]
Hence $(\partial_a^2+\partial_b^2+(2s/b)\partial_b)E(s)=0$
by continuation from absolute convergence. Differentiate in $s$ at
$s=-m$. The extra term $(2/b)\partial_bE(-m)$ is zero, because
$E(-m)$ vanishes identically as a function of the parameters. This
proves \eqref{qspec:eq:PDE}. The mass ladder yields
\eqref{qspec:eq:PDE-adjacent} for $m\ge2$.
\end{proof}
For $m=0$ the logarithm in \eqref{qspec:eq:L0} is harmonic on the open
strip $b>0$. The higher weights satisfy its explicit weighted version,
rather than an unrelated differential equation at each polylogarithm order.

\subsection{Products of several quadratic factors}
The parity mechanism extends beyond a single conjugate pair. The following
finite reduction uses no multivariable special functions at the first
spectral derivative.

\begin{theorem}[Several masses with an arbitrary even polynomial weight]
\label{qspec:thm:several-masses}
Let $P$ be a polynomial of degree at most $D$, let $r\ge1$, and let
$b_1,\ldots,b_r>0$. Define, initially in its absolute-convergence
half-plane,
\begin{equation}\label{qspec:eq:several-masses}
\mathcal Z_{P,\mathbf b}(s;a)=
\sum_{n\in\Z}P((n+a)^2)
\prod_{j=1}^r\bigl((n+a)^2+b_j^2\bigr)^{-s}.
\end{equation}
Its continuation is holomorphic at zero, and
\begin{equation}\label{qspec:eq:several-masses-finite}
\mathcal Z_{P,\mathbf b}'(0;a)
=\sum_{j=1}^r\sum_{\ell=0}^D
\frac{P^{(\ell)}(-b_j^2)}{\ell!}L_\ell(a,b_j).
\end{equation}
Combining this formula with \eqref{qspec:eq:L-polylog} gives a finite
polylogarithm reduction for every such spectrum, including repeated masses.
\end{theorem}
\begin{proof}
First work on a positive tail $x=n+a$ with all $b_j<x$. After factoring
out $x^{-2rs}$, the remaining factor is
\[
\exp\left(-s\sum_{j=1}^r\log(1+b_j^2/x^2)\right).
\]
It has only even powers of $x^{-1}$. Multiplication by $P(x^2)$ still
leaves only even powers of $x$. A possible Hurwitz pole at $s=0$
would require the resulting power to be $x^{-1}$, which is impossible.
Thus there are no resonant pole terms at zero in this common-scale
expansion.

Compare the product spectrum with $r^{-1}$ times the sum of the
individual weighted spectra evaluated at $rs$. Their constant and
linear terms in $s$ agree identically before summation; the remaining
coefficients are $O(s^2)$. Since there is no pole at zero, their first
derivatives agree after continuation as well. Remove a finite head
to justify the expansion for arbitrary masses. The ordinary head
logarithms add, so it contributes no discrepancy. The same argument
applies to the negative lattice tail, where evenness changes its
weight to $P(y_n^2)$ without a sign.

Finally the polynomial Taylor identity
\[
P(x^2)=\sum_{\ell=0}^D\frac{P^{(\ell)}(-b_j^2)}{\ell!}
 (x^2+b_j^2)^\ell
\]
expresses the $j$th individual weighted derivative as the inner finite
sum in \eqref{qspec:eq:several-masses-finite}. This proves the formula.
\end{proof}

For instance, $P(v)=v+c$ and two masses give
\[
\mathcal Z_{P,(b_1,b_2)}'(0;a)
=\sum_{j=1}^2\bigl(L_1(a,b_j)+(c-b_j^2)L_0(a,b_j)\bigr).
\]
Each summand has only elementary logarithms, $\Li_2$, and $\Li_3$.
The theorem is restricted to the first spectral derivative. Products
of logarithms in higher derivatives are not removed by this argument.

\subsection{Positive-integer resolvents: polygamma and rational polylogs}
At positive integers the same half-integer Bessel polynomial produces
polylogarithms of nonpositive order.

\begin{proposition}\label{qspec:prop:positive}
For $N\ge1$,
\begin{align}
E(N;a,b)={}&\frac{\sqrt\pi\Gamma(N-1/2)}{\Gamma(N)}b^{1-2N}\notag\\
&+\frac{2\pi^N b^{-N}}{\Gamma(N)}
\sum_{j=0}^{N-1}\frac{(N-1+j)!}
 {j!(N-1-j)!(4\pi b)^j}\Rea\Li_{j+1-N}(z).
\label{qspec:eq:positive-polylog}
\end{align}
An independent finite polygamma expression is
\begin{equation}\label{qspec:eq:positive-polygamma}
E(N;a,b)=2\Rea\sum_{r=1}^N A_{N,r}(b)
\bigl(\Xi_r(a+ib)+\Xi_r(1-a+ib)\bigr),
\end{equation}
where
\begin{align}
A_{N,r}(b)&=(-1)^{N-r}\binom{2N-r-1}{N-r}(-2ib)^{-(2N-r)},\label{qspec:eq:Ar}\\
\Xi_1(c)&=-\psi(c),\qquad
\Xi_r(c)=\frac{(-1)^r}{(r-1)!}\psi^{(r-1)}(c)\quad(r\ge2).
\label{qspec:eq:Xi}
\end{align}
\end{proposition}
\begin{proof}
Insert \eqref{qspec:eq:half-integer-K}, with $m=N-1$, into
\eqref{qspec:eq:poisson}. The remaining powers of $n$ give
\eqref{qspec:eq:positive-polylog}.

For the other formula, expand the rational function at $x=-ib$:
\[
\frac1{(x^2+b^2)^N}
=\sum_{r=1}^N\left\{\frac{A_{N,r}(b)}{(x+ib)^r}
+\frac{\overline{A_{N,r}(b)}}{(x-ib)^r}\right\}.
\]
The coefficient follows by expanding $(y-2ib)^{-N}$ at $y=x+ib=0$.
For $r\ge2$ summation gives $\zeta(r,c)=\Xi_r(c)$. The $r=1$
terms must be paired before summation: their coefficients add to zero,
and a common finite cutoff gives finite part $-\psi(c)$. No separate
value for a divergent harmonic series is assigned. Sum over $x_n$
and $y_n$ to obtain \eqref{qspec:eq:positive-polygamma}.
\end{proof}
For $N=1$ these formulas reduce to the classical resolvent
\begin{equation}\label{qspec:eq:resolvent}
E(1;a,b)=\frac\pi b
\frac{\sinh(2\pi b)}{\cosh(2\pi b)-\cos(2\pi a)}.
\end{equation}
For larger $N$, all polylogs in
\eqref{qspec:eq:positive-polylog} are rational functions of $z$, obtained
by repeatedly applying $z\partial_z$ to $z/(1-z)$.

\subsection{Gamma products and every mass primitive of their logarithm}
Define the real quantity
\begin{equation}\label{qspec:eq:Rgamma}
R(a,b)=\log\frac{\Gamma(a)^2\Gamma(1-a)^2}
 {\Gamma(a+ib)\Gamma(a-ib)\Gamma(1-a+ib)\Gamma(1-a-ib)}.
\end{equation}
All conjugate products in the denominator are positive real numbers.

\begin{proposition}[A Gamma-product dilogarithm integral]
\label{qspec:prop:gamma-integrals}
The classical product and reflection identities give
\begin{equation}\label{qspec:eq:Rhyperbolic}
R(a,b)=\log\frac{\cosh(2\pi b)-\cos(2\pi a)}{1-\cos(2\pi a)}.
\end{equation}
In particular, with $z_0=e^{2\pi ia}$,
\begin{equation}\label{qspec:eq:gamma-Li2}
\int_0^b R(a,t)\dd t
=\pi b^2-b\log(4\sin^2\pi a)
+\frac1\pi\Rea\bigl(\Li_2(z)-\Li_2(z_0)\bigr).
\end{equation}
More generally, for every $r\ge1$,
\begin{align}
&\frac1{(r-1)!}\int_0^b(b-t)^{r-1}R(a,t)\dd t\notag\\
&\quad=\frac{2\pi b^{r+1}}{(r+1)!}
-\frac{b^r}{r!}\log(4\sin^2\pi a)\notag\\
&\qquad\ -2\Rea\left\{\frac1{(-2\pi)^r}
\left(\Li_{r+1}(z)-\sum_{j=0}^{r-1}\frac{(-2\pi b)^j}{j!}
\Li_{r+1-j}(z_0)\right)\right\}.
\label{qspec:eq:gamma-all-primitives}
\end{align}
\end{proposition}
\begin{proof}
The Gamma product, with its conjugate linear terms cancelling, gives
\[
\sum_{n\ge0}\log\left(1+\frac{b^2}{(n+a)^2}\right)
=2\log\Gamma(a)-\log\Gamma(a+ib)-\log\Gamma(a-ib).
\]
Add the formula at $1-a$ and use Gamma reflection \cite{gamma}; equivalently
compare \eqref{qspec:eq:L0} at $b$ and at zero. The result is
\eqref{qspec:eq:Rhyperbolic}. With the logarithmic branches already fixed,
it can be written
\[
R(a,b)=2\pi b-\log(4\sin^2\pi a)-2\Rea\Li_1(e^{-2\pi b}z_0).
\]
One integration using $\partial_b\Li_2(z)=-2\pi\Li_1(z)$ proves
\eqref{qspec:eq:gamma-Li2}. For repeated integration, an $r$th primitive
of $\Li_1(z)$ is $(-2\pi)^{-r}\Li_{r+1}(z)$. Subtract its Taylor
polynomial through order $r-1$ at zero to enforce all anchored initial
conditions. This gives \eqref{qspec:eq:gamma-all-primitives}. Unit-circle
values appearing in that Taylor polynomial have order at least two and
are absolutely convergent.
\end{proof}

% END sections/06-polylogarithms

% BEGIN sections/07-audit-research
\section{Source audit, verification, and further research}
\label{qspec:sec:audit}

\subsection{What was inspected}
The repository anchor is \qpin. Its commit timestamp is
11 October 2026 at 05:06:38 UTC, which is 10 October 2026 in the user's
America/Los\_Angeles time zone. The local date in this delivery is therefore
10 October 2026.

The source reading included the manuscript's main file and preface, its
research-status chapter, the beginning and central conventions of the
integration chapter, the displayed $S_6$ target and formal-obstruction
framework in the full-Cayley continuation, the arithmetic spectral-jets
section on periodic weights and poles, and the all-rational formal
classification section of the Herglotz report. The intake README and
incoming directory catalogue were also read \cite{repo,arithmetic,classification,intake}.
This was a targeted reading, not a full audit of every included manuscript
chapter, report, proof prerequisite, or source program.

Six incoming archive names were visible at the pin: the Bernoulli--Tornheim--
digamma, cubic-tail-descent, determinantal-harmonic-resonance,
differential-polygamma-and-Lerch-jets, multishift-reflection, and
polygamma-tail-identities packages. Only their catalogue metadata was
inspected; their interiors were not read. Accordingly, no theorem,
conjecture, omission, or error is attributed here to any of those archive
interiors. The machine-readable provenance file and source-audit note
preserve this distinction.

\subsection{A concrete editorial synchronization correction}
The inspected research chapter still describes the general rational
$J(p/q)$ symbol classification as a further question. At the same pin,
the report \qfile{herglotz-rational-classification} already contains a
section entitled ``Complete formal classification at rational arguments''
with a theorem and its argument \cite{classification}.

Its statement concerns coprime positive $p,q$. In mixed parity, with
$e$ the even entry and $o$ the odd entry, formal rational-dilogarithmic
reducibility and formal logarithmic reducibility are equivalent to
\begin{equation}\label{qspec:eq:existing-J-status}
e^2\equiv\pm1\pmod o,\qquad o^2\equiv1\pmod{2e}.
\end{equation}
When both entries are odd, the stated rational-dilogarithmic cases have
$p,q\in\{1,3,5\}$, and the only logarithmic case is $p=q=1$.
This classification belongs to the existing report; it is not a result
of the present article. Its prerequisite symbol-theoretic lemmas were
not all re-audited here.

The proposed correction is therefore editorial: replace the stale open
classification question by a pointer to that existing theorem, subject
to the repository's normal integration review, while retaining the
separate tasks of constructing short explicit evaluations and understanding
numerical period relations. Failure of a formal reduction criterion is
not numerical transcendence or period independence. The delivered
\qfile{integration/editorial_update.md} contains precise replacement wording.

No new counterexample to an analytic formula in the inspected source
material was established in this session. The following mathematical
safeguards concern the present derivations; they are not allegations
about unread reports.

\subsection{Five distinctions required for a correct implementation}
\begin{enumerate}
\item \textbf{The doubled spectral variable.} The pole of
$\zeta(1+2\epsilon,a)$ is $1/(2\epsilon)$, and its $m$th regular
jet is $(-2)^m\gamma_m(a)$. Omitting the factor two breaks both
reflection channels and the anomaly residue calculation.
\item \textbf{A finite part is not a Taylor coefficient before pole removal.}
The coefficient $C_m$ is
$\Bcal_{m+1}(\ell)/(m+1)$, not $\Bcal_m(\ell)$. The extra index comes
from dividing the analytic function $\epsilon C(\epsilon)$ by
$\epsilon$.
\item \textbf{The binomial radius is real.} Formula
\eqref{qspec:eq:harmonic-counterterms} cannot simply be evaluated outside
$b<\min(a,1-a)$. The finite-head formula supplies the exact replacement.
Both tail shifts increase by $M$.
\item \textbf{Complex-order holomorphy is not a real-part operation.}
The combinations in \eqref{qspec:eq:zt} are needed before differentiating
in a complex spectral order. At real orders they may then be evaluated
as real and imaginary parts.
\item \textbf{Pointwise logarithms do not automatically factor continued sums.}
The odd-weight examples \eqref{qspec:eq:anomaly-examples} are explicit
counterexamples to unchecked multiplicativity. Evenness guarantees
universal vanishing for this conjugate-pair polynomial problem; a
fortuitous zero at one mass does not imply evenness.
\end{enumerate}

\subsection{Executed verification and its evidentiary meaning}
The supplied code was copied to a separate working directory before the
recorded runs. The results are included in
\qfile{data/exact_results.json} and \qfile{data/numeric_results.json}.
The exact suite ran with SymPy 1.14.0. The numerical suite ran with
mpmath 1.3.0 at 45 working decimal digits.

\begin{center}
\begin{tabular}{@{}p{0.65\linewidth}r@{}}
\toprule
Exact finite regression family & Assertions\\
\midrule
Pochhammer jets and Bell polynomials & 140\\
Ordinary/half-shifted harmonic conversion & 120\\
Anomaly coefficients, their nonvanishing, and parity zeros & 36\\
Rational partial fractions & 6\\
Half-integer Bessel recurrence & 12\\
Mass-primitive coefficient recurrence & 85\\
Shift/mass differential equation & 88\\
Heat-residue normalization & 13\\
Anchored primitive pole cancellation & 8\\
Several-mass polynomial change of basis & 9\\
\midrule
Total & 517\\
\bottomrule
\end{tabular}
\end{center}

The numerical suite passed 82 independent comparisons. Eighteen test the
two Stieltjes channels through order three, including masses outside the
unshifted binomial disk. Twenty-four compare negative-point derivatives
using the differentiated Hurwitz series, finite complex-Hurwitz formula,
and finite polylogarithm formula. Sixteen test even and odd polynomial
anomalies; ten compare positive-integer polygamma and polylogarithm
expressions. The remaining fourteen test general Poisson continuation,
Gamma ratios and integrals, parameter primitives, both polylogarithm
integrals, and the several-mass theorem, including repeated masses.

The largest recorded scaled residual was
$4.618\times10^{-44}$, where the scale is
$\max(1,|\mathrm{lhs}|,|\mathrm{rhs}|)$. The acceptance threshold was
$10^{-33}$. These are \emph{floating-point diagnostics}, not rigorous
interval enclosures, tail certificates, or proofs of equality. The
truncation choices are conservative numerical choices, not a delivered
validated error estimator. The exact finite suite proves its particular
algebraic assertions, not the analytic continuation or all-index theorems.
Those depend on the ordinary proofs given in the article.

\subsection{Integration recommendation and claim boundaries}
A suitable destination is a new report under
\qfile{Analysis/Polylogarithms/docs/reports/quadratic-spectrum-stieltjes/}.
The compact integration fragment uses prefixed labels and can be adapted
into the Hurwitz--Stieltjes and integration chapters after normal editorial
review. No automatic patch, remote commit, overwrite of existing research,
or change of conjecture status is included.

The key claims to preserve are: all-index exact identities in both
reflection channels; all-positive-mass convergence through a finite head;
a universal, not pointwise, parity characterization; finite first-derivative
reductions, not arbitrary higher-derivative factorization; and explicitly
anchored primitives. The classical inputs and the new derivations should
remain distinguishable in the merged text.

\subsection{Further research questions and concrete next targets}
\begin{enumerate}
\item \textbf{Higher spectral jets at negative integers.}
For $r\ge2$, determine the smallest explicitly stated function algebra
containing $\partial_s^rE(-m;a,b)$. The reciprocal-Gamma zero leaves
Bessel order derivatives of order at most $r-1$. Can the second jet be
reduced to a finite family of polylogarithm order derivatives plus
explicit logarithmic terms, or is an additional integral coordinate
required? A failure of a chosen formal calculus must be distinguished
from numerical nonreducibility.

\item \textbf{Higher-jet factorization defects.}
For a fixed polynomial $W$, study
\[
\Delta_r(W)=2^{r-1}\bigl(F_{W,+}^{(r)}(0)+F_{W,-}^{(r)}(0)\bigr)
-Z_W^{(r)}(0).
\]
At $r=1$ Theorem~\ref{qspec:thm:anomaly} gives an exact local polynomial.
At $r=2$ mixed logarithms introduce a genuinely nonlocal part. Separate
that part from the residue polynomial and determine what parity still
forces to vanish. This is an extension problem, not an assumption that
the first-derivative criterion persists unchanged.

\item \textbf{Noninteger polynomial degree.}
Replace $W(x)$ by $x^\rho$ on the positive spectrum. Determine the exact
continuation as $\rho$ crosses the resonant odd integers and formulate a
parameter-dependent correction that specializes to
\eqref{qspec:eq:anomaly-harmonic}. The finite polynomial residue in this
article provides a set of exact interpolation constraints.

\item \textbf{Arithmetic projection of both Stieltjes channels.}
At $a=p/q$, apply finite character projectors to the master identities
and their mass-transfer laws. Which combinations admit a shorter
Gamma/Dirichlet-$L$ jet description, and which retain genuinely useful
damped polylogarithm integrals? The first useful target is a complete
small-denominator dictionary with the even and odd character components
kept separate.

\item \textbf{Exact removal of the auxiliary mass.}
The identities are independent of $b$ after their counterterms are
combined. Develop an exact coefficient calculus for the cancellations
as $b\downarrow0$, rather than simply taking a numerically unstable
limit of separate terms. Comparing the resulting coefficients at two
masses may generate new harmonic sums involving Hurwitz order derivatives.

\item \textbf{Mixed shift and mass primitives.}
Combine \eqref{qspec:eq:all-primitives},
\eqref{qspec:eq:mass-primitives}, and \eqref{qspec:eq:PDE} to classify
which rectangular integrals in $(a,b^2)$ have finite polylogarithm and
Hurwitz-jet endpoints. Anchoring every primitive should be part of the
statement, not a later adjustment by unspecified polynomials.

\item \textbf{Several factors beyond the first spectral derivative.}
Theorem~\ref{qspec:thm:several-masses} has a finite answer even with
repeated masses. Determine the second derivative for two masses and
its confluent limit as they coincide. The target is an exact formula
for the logarithmic cross term, with no subtraction of divergent sums
left implicit.

\item \textbf{Formalization and validated evaluation.}
The Bell recurrences, harmonic coefficient identity, rational partial
fractions, and finite mass/PDE reductions are suitable finite targets
for a proof assistant. Separately, an interval implementation of the
Bessel order integrals and finite-head remainder would turn the present
floating-point diagnostics into validated evaluations. Neither has
been supplied here.
\end{enumerate}

\section*{Conclusion}
\addcontentsline{toc}{section}{Conclusion}
The positive quadratic deformation does more than provide another series
for a Stieltjes constant. Its even and odd spectra form a complete
reflection pair, its harmonic counterterms admit a globally valid
finite-head formulation, and its negative spectral zeros collapse first
derivatives to finite polylogarithm expressions. Polynomial parity
explains exactly when conjugate factorization is legitimate. The same
finite family supports anchored mass primitives, a shift/mass differential
equation, and products of several quadratic factors. These are exact
structural identities with ordinary proofs; the supplied computations
check their implementation without replacing those proofs.

% END sections/07-audit-research

% BEGIN references
\begin{thebibliography}{99}
\bibitem{repo}
V. Reshetnikov, \emph{ProveIt: Polylogarithms manuscript}, source snapshot
\qpin. In particular the main manuscript, integration chapter,
research-status chapter, and full-Cayley obstruction continuation.
\url{https://github.com/VladimirReshetnikov/ProveIt/tree/446dca9bd65c20e409d1d2561d4ffad408103e2a/Analysis/Polylogarithms/docs/manuscript}.

\bibitem{arithmetic}
\emph{Arithmetic spectral jets}, ProveIt research continuation,
section ``Periodic weights and several distinct arithmetic poles'', same pin.
\qfile{Analysis/Polylogarithms/docs/reports/stieltjes-harmonic-resonance/continuations/arithmetic-spectral-jets/sections/04-arithmetic.tex}.
The inspected section credits the classical simple-pole anomaly mechanism.

\bibitem{classification}
\emph{Herglotz rational classification}, ProveIt research report,
section ``Complete formal classification at rational arguments'', same pin.
\qfile{Analysis/Polylogarithms/docs/reports/herglotz-rational-classification/sections/04-classification.tex}.
This is the source of the editorial status correction, not a new
classification claimed in this article.

\bibitem{intake}
ProveIt, \emph{Incoming reports}, intake instructions and archive catalogue,
same pin. \qfile{docs/incoming/README.md} and \qfile{docs/incoming/}.
Only archive metadata, not archive interiors, was inspected.

\bibitem{hurwitz}
NIST Digital Library of Mathematical Functions, \S25.11,
\emph{Hurwitz zeta function}.
\url{https://dlmf.nist.gov/25.11}. Consulted 10 October 2026.

\bibitem{polylog}
NIST Digital Library of Mathematical Functions, \S25.12,
\emph{Polylogarithms}.
\url{https://dlmf.nist.gov/25.12}. Consulted 10 October 2026.

\bibitem{bessel}
NIST Digital Library of Mathematical Functions, \S10.32,
\emph{Integral representations of modified Bessel functions}, especially
(10.32.9) and (10.32.10).
\url{https://dlmf.nist.gov/10.32}. Consulted 10 October 2026.

\bibitem{halfK}
NIST Digital Library of Mathematical Functions, \S10.39,
\emph{Relations to other functions}, especially the half-integer base case;
the higher finite formula used here is also derived by recurrence.
\url{https://dlmf.nist.gov/10.39}. Consulted 10 October 2026.

\bibitem{theta}
NIST Digital Library of Mathematical Functions, \S20.7,
\emph{Theta-function identities and transformations}.
\url{https://dlmf.nist.gov/20.7}. Consulted 10 October 2026.

\bibitem{gamma}
NIST Digital Library of Mathematical Functions, \S5.5,
\emph{Functional relations}, and \S5.15, \emph{Polygamma functions}.
\url{https://dlmf.nist.gov/5.5}; \url{https://dlmf.nist.gov/5.15}.
Consulted 10 October 2026.

\bibitem{dowker}
J. S. Dowker, \emph{Calculation of the multiplicative anomaly},
arXiv:1412.0549, 2014; corrected version 3, 24 September 2023.
\url{https://arxiv.org/abs/1412.0549}.
Cited for the classical subject and scope; the precise defect and sign
convention used in this article are proved independently.
\end{thebibliography}

% END references

\end{document}
